\textbf{COLUNI 2023 -- Questão 36}\par\medskip “O Estado Novo se constituiu em decorrência de uma política de massas que se foi definindo no Brasil a partir da Revolução de 1930, com a ascensão de Getúlio Vargas ao poder. Esse tipo de política, voltada para as classes populares, desenvolveu-se no período entre as guerras a partir das críticas ao sistema liberal, considerado incapaz de solucionar os problemas sociais. [...] Uma das soluções propostas era a do controle social através da presença de um Estado forte comandado por um líder carismático, capaz de conduzir as massas no caminho da ordem”.

\par

CAPELATO, Maria Helena. O Estado Novo: o que trouxe de novo?. In: FERREIRA, Jorge e DELGADO, Lucilia de Almeida. \textbf{O Brasil Republicano:} O tempo do nacional-estatismo. v. 2. Rio de Janeiro: Civilização Brasileira, 2007. p. 109.

\par

Tendo como base o trecho acima e os conhecimentos históricos sobre o Estado Novo, assinale a alternativa INCORRETA:\begin{enumerate}
\item Uma das principais mudanças, no Estado Novo, em relação ao cenário da Primeira República, foi a
      transformação da questão social como questão de Estado e não questão de polícia.
\item A cerimônia cívica com a queima das bandeiras estaduais representou um ato simbólico da vitória do
      projeto centralizador do Estado Novo.
\item O poder político no Estado Novo passou a se apoiar nas forças populares e liberais da nação, com a
      liderança de grupos urbanos na vida política do país.
\item A propaganda política e a repressão aos opositores foram dois grandes pilares de sustentação do Estado
      Novo varguista e formas de se assegurar sua legitimidade.\end{enumerate}
